\documentclass[11pt]{article}
\usepackage[utf8]{inputenc}\usepackage[T1]{fontenc}\usepackage{mathptmx}
\usepackage[english]{babel}
\usepackage[a4paper,margin=2.2cm]{geometry}
\setlength{\parindent}{0pt}\setlength{\parskip}{7pt}\linespread{1.05}
\pagestyle{empty}
\begin{document}
\hfill Lima, 12 November 2026

\begin{center}\large\bfseries Cover Letter\end{center}

Dear Editor,

We are pleased to submit our manuscript entitled “Machine Learning Model for Predicting the Risk of Childhood Stunting in Peru” for consideration as an original research article in \emph{Nutrition}. The authors are Abel Is-Boset Guevara-Huasco, Zulin Pamela Vallejos-Cotrina, Verónica Lisseth Vergara-Rojas, Ferdinand Edgardo Pineda-Ancco and Nemias Saboya-Rios.

The study develops and internally validates a model to predict stunting in 304,271 children under five, using the Peruvian Demographic and Family Health Survey (ENDES) 2007–2025 and 45 predictors from five domains: the child, the mother, the household and territory, water and sanitation, and housing. Six supervised algorithms were compared under stratified five-fold cross-validation with survey weights. LightGBM was selected, with an out-of-fold AUC-ROC of 0.829 (95\% CI 0.827–0.832) and a sensitivity of 76.4\%. SHAP values identified household poverty, birth weight and altitude as the leading predictors and showed how their weight differs across Peru’s 25 departments.

The model was implemented as a web-based decision-support application, which a panel of four health professionals rated favorably (median overall score of 9 out of 10). The code, the processed dataset and the trained model are openly available in Zenodo to support reproducibility.

We believe the manuscript will interest the readers of \emph{Nutrition}, which has published several machine learning studies on child undernutrition in Bangladesh, Egypt, India, the Philippines and sub-Saharan Africa. Our work extends this line of research to Latin America, where it remains scarce, and adds a territorial interpretation and a practical tool aimed at public health institutions.

We confirm that this manuscript has not been published elsewhere and is not under consideration by any other journal. All authors have approved the manuscript and agree with its submission to \emph{Nutrition}. The study was approved by the Ethics Committee of the School of Systems Engineering, Universidad Peruana Unión (No. CEI-UPEU-FIA-2026-24), and the authors declare no competing interests.

Thank you for your time and consideration.

Sincerely,

\vspace{1.2cm}
{\setlength{\parskip}{0pt}\small \textbf{Nemias Saboya-Rios, MSc}\\Corresponding author\\School of Systems Engineering, Faculty of Engineering and Architecture\\Universidad Peruana Unión, Carretera Central Km 19.5, Ñaña, Lurigancho-Chosica, Lima, Peru\\saboya@upeu.edu.pe · ORCID 0000-0002-7166-2197\par}
\end{document}
